\section*{Chapter \ref{ch:b2:polymer-synthesis} --- Polymers: Synthesis, Structure and Properties}

\begin{solution}{exo:b2:polymer-synthesis:1}
Polyethene \ce{-[CH2-CH2]-}; polypropene \ce{-[CH2-CH(CH3)]-}; poly(vinyl chloride)
\ce{-[CH2-CHCl]-}; polystyrene \ce{-[CH2-CH(C6H5)]-}; nylon-6,6 \ce{-[NH(CH2)6NH-CO(CH2)4CO]-}.
\end{solution}

\begin{solution}{exo:b2:polymer-synthesis:2}
PET: step growth (polyester, water released). Polystyrene: chain growth. Nylon-6,6: step growth.
Poly(methyl methacrylate): chain growth (a \ce{C=C} \omterm{def:b2:polymer-synthesis:polymer}{monomer}). Poly(lactic acid) from the hydroxy acid:
step growth, an \ce{A-B} \omterm{def:b2:polymer-synthesis:polymer}{monomer}.
\end{solution}

\begin{solution}{exo:b2:polymer-synthesis:3}
$\bar M_n = 1000 \times \qty{104.15}{g/mol} \approx \qty{104}{kg/mol}$ (styrene \ce{C8H8},
\qty{104.152}{g/mol}).
% ledger: aw:C, aw:H
\end{solution}

\begin{solution}{exo:b2:polymer-synthesis:4}
\omterm{def:b2:polymer-synthesis:tacticity}{Syndiotactic}. Yes: a regular chain can pack into crystallites.
\end{solution}

\begin{solution}{exo:b2:polymer-synthesis:5}
Mass fractions 0.5 and 0.5. $1/\bar M_n = 0.5/10 + 0.5/100 = 0.055$, $\bar M_n \approx
\qty{18.2}{kg/mol}$; $\bar M_w = 0.5 \times 10 + 0.5 \times 100 = \qty{55}{kg/mol}$; $\text{\DH} \approx
3.0$.
\end{solution}

\begin{solution}{exo:b2:polymer-synthesis:6}
$1/(1-0.98) = 50$; $1/(1-0.995) = 200$.
\end{solution}

\begin{solution}{exo:b2:polymer-synthesis:7}
$r = 1/1.01$; at $p = 1$, $\bar X_n = (1+r)/(1-r) = (1.01 + 1)/(1.01 - 1) = 201$.
\end{solution}

\begin{solution}{exo:b2:polymer-synthesis:8}
$R_p \propto [\mathrm I]^{1/2}$: multiplied by $\sqrt2 \approx 1.41$. $\nu \propto [\mathrm I]^{-1/2}$:
divided by $\sqrt 2$, about 0.71 times its former value.
\end{solution}

\begin{solution}{exo:b2:polymer-synthesis:9}
Pure PVC has its glass transition above room temperature: a rigid glass. The small plasticiser
molecules sit between the chains, separate them and let segments move at lower temperature: $T_g$ falls
below room temperature, and the material is rubbery and flexible in use.
\end{solution}

\begin{solution}{exo:b2:polymer-synthesis:10}
$\dfrac{\mathrm d}{\mathrm dx}\bigl(x p^{\,x-1}\bigr) = p^{\,x-1}(1 + x\ln p)$, zero for $x = -1/\ln p$
(a maximum: the sign goes from $+$ to $-$). For $p$ close to 1, $\ln p \approx -(1-p)$, so $x \approx
1/(1-p) = \bar X_n$.
\end{solution}

\begin{solution}{exo:b2:polymer-synthesis:11}
Stage 1, \omterm{def:b2:acyl-substitution:transesterification}{transesterification}: \ce{C6H4(COOCH3)2 + 2HOCH2CH2OH -> C6H4(COOCH2CH2OH)2 + 2CH3OH}, methanol
distilled off. Stage 2: each diester molecule carries two \ce{CH2CH2OH} ends; one end attacks an ester
of another molecule and releases a molecule of ethane-1,2-diol, which is pumped off. The \omterm{def:b2:polymer-synthesis:polymer}{monomer}
carries both kinds of group in the right ratio, as an \ce{A-B} \omterm{def:b2:polymer-synthesis:polymer}{monomer} does: the balance $r = 1$ is
built in, and removing the diol drives $p$ towards 1.
\end{solution}

\begin{solution}{exo:b2:polymer-synthesis:12}
Equimolar feed: ratio $= (2.0 + 1)/(1 + 0.5) = 2$: twice as many units of \omterm{def:b2:polymer-synthesis:polymer}{monomer} 1. With $r_1 = r_2 =
0$, ratio $= [\mathrm M_1][\mathrm M_2]/([\mathrm M_2][\mathrm M_1]) = 1$ whatever the feed: each
radical adds only the other \omterm{def:b2:polymer-synthesis:polymer}{monomer}, an alternating \omterm{def:b2:polymer-synthesis:copolymer}{copolymer}.
\end{solution}

\begin{solution}{pb:b2:polymer-synthesis:1}
\textbf{1.} \ce{H2N(CH2)6NH2} and \ce{HOOC(CH2)4COOH}.
\textbf{2.} One water per link:
\[\mathrm{{-}COOH + H_2N{-} \longrightarrow {-}CO{-}NH{-} + H_2O}.\]
\textbf{3.} The salt contains exactly one diamine per diacid; weighing two separate liquids or solids
would never reach the exactness the Carothers equation demands.
\textbf{4.} \ce{-[NH(CH2)6NH-CO(CH2)4CO]-}, \ce{C12H22N2O2}, \qty{226.32}{g/mol}.
\textbf{5.} The \omterm{def:b2:polymer-synthesis:polymer}{repeat unit} contains two monomer-derived units: $M_0 = 226.32/2 = \qty{113.16}{g/mol}$.
\textbf{6.} The carbons of the diamine (6) and of the diacid (6).
\textbf{7.} $\bar X_n = 50$, $\bar M_n = 50 \times 113.16 \approx \qty{5.66}{kg/mol}$.
\textbf{8.} $\bar X_n = 100$, $\bar M_n \approx \qty{11.3}{kg/mol}$.
\textbf{9.} With 99~\% of the groups reacted, the average chain has only 100 units, too short for a
strong fibre; each further step towards complete \omterm{def:b2:continuous-reactors:conversion}{conversion} doubles or triples the chain length.
\textbf{10.} $\bar X_n = 10\,000/113.16 = 88.4$; $p = 1 - 1/88.4 \approx 0.989$.
\textbf{11.} By removing water: the melt is heated well above its melting point, finally under reduced
pressure, so that the equilibrium keeps moving towards amide.
\textbf{12.} A monofunctional impurity caps chain ends, and any impurity upsets the 1\,:\,1 balance.
\textbf{13.} $r = 1/1.01 \approx 0.990$.
\textbf{14.} $\bar X_n = (1+r)/(1-r) = 201$; $\bar M_n \approx 201 \times 113.16 \approx
\qty{22.7}{kg/mol}$.
\textbf{15.} $\bar X_n = 1.9901/(1.9901 - 2 \times 0.9901 \times 0.99) \approx 1.9901/0.0297 \approx
67$.
\textbf{16.} Amine groups: when the acid groups are used up, every chain ends in \ce{NH2}.
\textbf{17.} It turns an amine end into an amide that cannot react further: it caps chains, limiting
$\bar M_n$ and stopping further growth when the \omterm{def:b2:polymer-synthesis:polymer}{polymer} is remelted.
\textbf{18.} To keep a reproducible melt viscosity for spinning and moulding, which would otherwise
drift as chains keep reacting in the melt.
\textbf{19.} $\text{\DH} = 1 + p = 1.99$.
\textbf{20.} $\bar M_w = 1.99 \times 11.3 \approx \qty{22.5}{kg/mol}$.
\textbf{21.} Number fraction $n_1 = 1 - p = 0.01$, 1~\% of the molecules; mass fraction $w_1 = (1-p)^2 =
10^{-4}$, 0.01~\% of the mass.
\textbf{22.} Near $x = -1/\ln 0.99 \approx 99.5$, that is near $\bar X_n = 100$ (exercise 10).
\textbf{23.} $1000/226.32 = \qty{4.42}{mol}$ of \omterm{def:b2:polymer-synthesis:polymer}{repeat units}, two amide links each: \qty{8.84}{mol} of
water, $8.84 \times 18.015 \approx \qty{159}{g}$.
\textbf{24.} Chains too short give a weak fibre; chains too long give a melt too viscous to be pushed
through the fine holes of the spinneret.
\textbf{25.} $\bar X_n = 20\,000/113.16 = 176.7$ and $p = 1 - 1/176.7 = \boldsymbol{\approx 0.994}$.
\end{solution}
